\documentclass{article}



\begin{document}
	$ $
	\vspace{9cm}
	\begin{center}
		\textbf{\Large{Caderno de Cálculo Numérico}}
	
		Leandro Júnior
	\end{center}
	\newpage
	\textbf{Noções Básicas sobre ERROS}
	
	\textbf{Data:} 22/09/2026
	
	\hrule 
	\vspace{1cm}
		
	\textbf{\large{1. Representação dos Números}}
	\vspace{0.5cm}
	
	\textbf{Exemplo 1}
	
	
	Calcular a área de um circunferência com $100$ m  de raio:
	
	\begin{enumerate}
		\item $A = 31400$ m$^2$
		\item $A = 31416$ m$^2$
		\item $A = 31415,92654$ m$^2$
	\end{enumerate}
		
		Por que a área pôde assumir mais de um valor?
		
		A área de um circunferência de raio $R$ é calculada através da expressão $\pi R^2$. Ou seja, depende do número irracional $\pi$, isso significa que depende de um número que é impossível representar com um número finito de casas decimais, quanto mais casas maior a precisão na representação. Em cada exemplo $\pi$ foi representado com diferentes precisões $3,14$ em (1); $3,1416$ em (2) e $3,14159654$ em (3). É claro que no terceiro caso a área $A$ se aproximou mais da área real que é impossível de ser representada exatamente por um número.
		\vspace{1cm}
		
		\textbf{1.1 Conversão entre números da base decimal para a base binária}
		\vspace{0.5cm}
		
		Um número como $(235)_{10}$ pode ser escrito dessa forma:
		\begin{center}
			\begin{tabular}{lll}
				$(235)_{10}$&$=$&$2\cdot 10^2 + 3\cdot 10^1 + 5 \cdot 10^0$
			\end{tabular}
		\end{center}
		
		Na realidade qualquer número de uma base $\beta$ $(a_ka_{k-1}\dots a_1a_0)_\beta$ com $0\leq a_i <\beta$ pode ser representado da seguinte maneira:
		
		\[a_k\beta^{k} + a_{k-1}\beta^{k-1} + \cdots + a_0\beta^0\]
		
		Sendo assim o número $(1011)_2$ é igual a:
		\[2^3 + 0\cdot 2^2 + 2^1 + 2^0\]
		
		Deixando o $2$ em evidência:
		
		\[2\cdot(2^2+1) + 1\]
		\[2\cdot(2\cdot 2 +1) + 1\]
		\[2\cdot(2\cdot(2\cdot1+0)+1)+1\]
		
		Resolvendo a expressão seguindo a ordem dos parênteses; 
		
		\[2\cdot(2\cdot(2+0)+1)+1\]
		\[2\cdot(4+1)+1\]
		\[2\cdot(5)+1\]
		\[10+1\]
		\[11\]
		
		Outro exemplo seria $(111)_2$:
		
		\[2^2+2^1+2^0\]
		\[2\cdot(2\cdot 1+1)+1\]
		\[2\cdot(3)+1\]
		\[6+1\]
		\[7\]
		
		Através desses exemplos é possível construir um processo de converter um número binário para um número decimal:
		
		Sendo o número binário $(b_kb_{k-1}b_{k-2}\dots b_1b_0)_2$ a representação dele em base decimal pode ser obtida da seguinte forma:
		
		\begin{center}
			\begin{tabular}{lll}
				$p_1$ &$=$& $b_k$\\
				$p_2$&$=$&$2\cdot p_1+b_{k-1}$\\
				$p_3$&$=$&$2\cdot p_2 + b_{k-2}$\\
				$\vdots$&$\vdots$&$\vdots$\\
				$p_{k+1}$&$=$&$2\cdot p_k + b_0 = N_{10} $
			\end{tabular}
		\end{center}
		
		
	
	
\end{document}